\documentclass[11pt]{article}
\usepackage[margin=1in]{geometry}
\usepackage{amsmath,amssymb}
\usepackage{booktabs}
\usepackage{hyperref}

\title{Scenario 1 --- Shrinking-Horizon Certainty-Equivalent MPC\\
       for Mobile Charging Station Dispatch\\[2pt]
       \large Formal Model}
\author{Approach 1 (Deterministic Certainty-Equivalent MPC)}
\date{}

\newcommand{\dt}{\Delta t}

\begin{document}
\maketitle

Formal companion to the code in \texttt{Shrinking\_Horizon/code/} (the window MILP is
built in \texttt{3\_MCSModel.jl}; the closed loop in \texttt{4\_MPCLoop.jl}). Plain-English
descriptions and the data schema are in \texttt{README.md}; a code audit is in
\texttt{constraints\_code\_vs\_model.txt}.

The controller is a \textbf{single-day, full-24\,h shrinking horizon}. Each 15-minute step
re-solves a MILP over the \emph{entire remaining day} $[k_0,\dots,n_{\text{day}}]$; only the
first interval is applied, and the window shrinks as the day advances. There is one
optimisation per step and one horizon: the overnight MCS recharge is scheduled \emph{inside}
this MILP (there is no separate phase). The activity powers are a \textbf{fixed} calibrated
estimate; the MILP plans on the mean, and the stochastic plant samples per-interval powers
from $\mathcal{N}(\mu,\sigma)$.

%======================================================================
\section{Sets and indices}
\begin{tabular}{ll}
$m \in \mathcal{M}$ & mobile charging stations (MCS) \\
$e \in \mathcal{E}$ & construction EVs (excavators) \\
$i,j \in \mathcal{N}$ & nodes; $\mathcal{N}_g$ grid, $\mathcal{N}_c$ sites \\
$a \in \mathcal{B}$ & activities $=\{\text{dig},\text{load},\text{trv},\text{idle}\}$ \\
$k$ & interval $1,\dots,n_{\text{day}}$; window $[k_0,n_{\text{day}}]$ \\
$\mathcal{T}_b$ & interval boundaries of the window \\
\end{tabular}

\medskip
$n_{\text{day}} = n_{\text{int}} = 96$ is the \textbf{full 24\,h} day at $\dt=0.25$\,h
(08:00 $\to$ 08:00 next day). $A_{i,e}=1$ iff CEV $e$ works at site $i$. Work availability
$R_{i,e,k}$ is $0$ outside the shift, but the MCS/CEVs may charge and the MCS may travel at
any interval.

%======================================================================
\section{Parameters (selected)}
\begin{tabular}{ll}
$\dt$ & interval length (h), $0.25$ \\
$\lambda^{\text{buy}}_{k},\lambda^{\text{CO2}}_{k}$ & price (\$/kWh) / carbon intensity \\
$\overline{SOE}_\bullet,\underline{SOE}_\bullet,SOE^{\text{ini}}_\bullet$ & energy cap / floor / start (MCS, CEV) \\
$CH_m,DCH_m,DCH^{\text{plug}}_m,C^{\text{plug}}_m,\eta_m,CH^{\text{CEV}}_e$ & MCS / CEV charge-discharge parameters \\
$p_a$ & activity power (kW); fixed calibrated mean, MILP plans on it \\
$H^{\text{dig}}_i,H^{\text{load}}_i$ & lumpsum dig / load requirement (hours) at site $i$ \\
$\tau_{i,j},k_{\text{trv}},R_{i,e,k}$ & travel time / energy / work availability \\
$\rho_{\text{miss}},\rho_{\text{lab}},\lambda^{\text{NC}},\lambda^{\text{OP}}$ & penalties \\
$\text{scale},w_{\text{pt}}=4,r$ & precedence ratio / travels-per-work / rest cap \\
\end{tabular}

%======================================================================
\section{Decision variables}
\textbf{Continuous ($\ge0$):} $P^{\text{ch}}_{m,i,k}$, $P^{\text{dch}}_{m,i,k}$,
$P^{\text{MCS-CEV}}_{m,i,e,k}$, $P^{\text{work}}_{i,e,k}$, totals
$P^{\text{ch,tot}}_{m,k},P^{\text{dch,tot}}_{m,k}$, travel energy $L_{m,i,j,k}$, state
$SOE^{\text{MCS}}_{m,t},SOE^{\text{CEV}}_{e,t}$, peaks $P^{\text{NC}},P^{\text{OP}}$, and the
single work-shortfall slack $s^{\text{miss}}_{i,a}$.

\smallskip
\textbf{Binary:} $u_{e,i,a,k}$, $\mu_{i,e,k}$, $\rho_{m,i,e,k}$, $z_{m,i,k}$,
$g^{\text{ch}}_{m,i,k}$, $x_{m,i,j,k}$, $y_{m,i,j,k}$, $\beta^{\text{arr}},\beta^{\text{dep}}$.

%======================================================================
\section{Objective (Eq.\ 1)}
\begin{align}
\min\ &\sum_{m,k}\lambda^{\text{buy}}_{k}P^{\text{ch,tot}}_{m,k}\dt
      +\sum_{m,k}\tfrac{c_{\text{ton}}}{1000}\lambda^{\text{CO2}}_{k}P^{\text{ch,tot}}_{m,k}\dt
      +\rho_{\text{miss}}\!\!\sum_{i\in\mathcal{N}_c,a}\!\! s^{\text{miss}}_{i,a}\notag\\
     &+\lambda^{\text{NC}}P^{\text{NC}}+\lambda^{\text{OP}}P^{\text{OP}}
      +\rho_{\text{lab}}\dt\!\!\sum_{m,i,j,k}\!\! y_{m,i,j,k}.
\end{align}

%======================================================================
\section{Constraints}

\paragraph{Power aggregation \& flow (HARD).}
$P^{\text{ch,tot}}_{m,k}=\sum_{i\in\mathcal{N}_g}P^{\text{ch}}_{m,i,k}$,
$P^{\text{dch,tot}}_{m,k}=\sum_{i\in\mathcal{N}_c}P^{\text{dch}}_{m,i,k}$;
$P^{\text{dch}}_{m,i,k}=\sum_e P^{\text{MCS-CEV}}_{m,i,e,k}\le DCH_m z_{m,i,k}$;
$P^{\text{MCS-CEV}}_{m,i,e,k}\le DCH^{\text{plug}}_m\rho_{m,i,e,k}$;
$\sum_m P^{\text{MCS-CEV}}_{m,i,e,k}\le CH^{\text{CEV}}_e\mu_{i,e,k}$.
Grid exclusivity: $P^{\text{ch}}_{m,i,k}\le CH_m g^{\text{ch}}_{m,i,k}$,
$g^{\text{ch}}\le z$, $\sum_m g^{\text{ch}}_{m,i,k}\le 1$.

\paragraph{State of energy (HARD).}
\begin{align}
SOE^{\text{MCS}}_{m,k+1}&=SOE^{\text{MCS}}_{m,k}+\eta_m P^{\text{ch,tot}}_{m,k}\dt-\tfrac{1}{\eta_m}P^{\text{dch,tot}}_{m,k}\dt-L^{\text{tot}}_{m,k},\\
SOE^{\text{CEV}}_{e,k+1}&=SOE^{\text{CEV}}_{e,k}+\textstyle\sum_{m,i}P^{\text{MCS-CEV}}_{m,i,e,k}\dt-\sum_i P^{\text{work}}_{i,e,k}\dt,\\
\underline{SOE}_\bullet&\le SOE_\bullet\le\overline{SOE}_\bullet .
\end{align}

\paragraph{Terminal energy targets (Eq.\ 8a / 8b, HARD).}
Applied at the final window boundary:
\begin{align}
SOE^{\text{MCS}}_{m,n_{\text{day}}+1}&= SOE^{\text{MCS,ini}}_m && \text{(8a, exact; overnight refill inside this MILP)},\\
SOE^{\text{CEV}}_{e,n_{\text{day}}+1}&\ge SOE^{\text{CEV,ini}}_e && \text{(8b, floor; overcharging allowed)}.
\end{align}
The CEV target is a \emph{floor} rather than Avik's exact equality: a CEV cannot discharge, so
under the stochastic plant an equality can become unrecoverable; the floor keeps the terminal
reachable while guaranteeing the fleet ends at least as charged as it began.

\paragraph{Plugging, presence \& routing (Eq.\ 10, HARD).}
$\rho\le A$, $\rho\le z$, $\sum_e \rho_{m,i,e,k}\le C^{\text{plug}}_m$. Presence partition
$\sum_i z_{m,i,k}+\sum_{i\ne j}y_{m,i,j,k}=1$; $\beta^{\text{dep}}_{m,i,k}=\sum_j x_{m,i,j,k}$;
$\beta^{\text{arr}}$ from finishing trips (or a carried-in arrival);
$\beta^{\text{arr}}-\beta^{\text{dep}}=z_k-z_{k-1}$; $\beta^{\text{arr}}+\beta^{\text{dep}}\le1$;
flow balance (generalised for the carried-in start). Trip on $i\!\to\!j$ lasts $\tau_{i,j}$ and
costs $L=k_{\text{trv}}\dt\,y$. \textbf{(10e)} MCS parked at a grid node by day-end:
$\sum_{i\in\mathcal{N}_g}z_{m,i,n_{\text{day}}}=1$.

\paragraph{Activity scheduling (Eq.\ 11, HARD).}
$\sum_a u_{e,i,a,k}=A_{i,e}$; $u\le A$; $\mu_{i,e,k}\le u_{e,i,\text{idle},k}$;
$P^{\text{work}}_{i,e,k}\le R_{i,e,k}A_{i,e}(1-\mu_{i,e,k})$;
$P^{\text{work}}_{i,e,k}=\sum_a p_a u_{e,i,a,k}$ (Eq.\ 5e), with $p_{\text{idle}}=0$.
(The code uses this 4-activity encoding; Avik's equivalent 3-activity form is
$\sum_a u\le1$, $\sum_a u+\mu\le1$.)

\paragraph{Work quota (Eq.\ 12c, SOFT) --- lumpsum.}
\begin{equation}
s^{\text{miss}}_{i,\text{dig}}\ \ge\ H^{\text{dig}}_i-\dt\!\!\sum_{e,k}\!u_{e,i,\text{dig},k},\qquad
s^{\text{miss}}_{i,\text{load}}\ \ge\ H^{\text{load}}_i-\dt\!\!\sum_{e,k}\!u_{e,i,\text{load},k}.
\end{equation}
(No hard upper cap: with a single day there is no incentive to overshoot.)

\paragraph{Precedence (Eq.\ 12d, HARD).}
Raw interval counts, seeded by realized work carried in:
$\sum_{e,\tau\le k}u_{e,i,\text{load},\tau}\le\text{scale}\sum_{e,\tau\le k}u_{e,i,\text{dig},\tau}$.

\paragraph{Rest rule (Eq.\ 12e, HARD).}
With $r=\lceil t_{\text{rest}}/\dt\rceil$ (=4),
$\sum_{k'=k}^{k+r}\sum_{a\in\{\text{dig,load,trv}\}}u_{e,i,a,k'}\le r$, \emph{seeded from the
applied Work/Break history} so a work-run cannot leak across the every-15-min re-solves.

\paragraph{Travel pacing (Eq.\ 13, HARD).}
With $w_{\text{pt}}=4$, for each $(i,e)$ and $k$, on cumulative travel $V$ vs cumulative
useful work $W$ (dig+load):
\begin{equation}
W(k)-w_{\text{pt}}\ \le\ w_{\text{pt}}\,V(k)\ \le\ W(k),
\end{equation}
seeded with the travel/work already applied in earlier windows.

\paragraph{Demand peaks (E1, HARD).}
$P^{\text{NC}}\ge$ carried-in peak and $\ge\sum_m P^{\text{ch,tot}}_{m,k}$ for all $k$;
$P^{\text{OP}}$ likewise on on-peak $k$ only.

%======================================================================
\section{Single horizon, fixed model}
There is no overnight ``Phase~2'': the terminal rules (8a/8b) close the MCS energy cycle inside
the one 24\,h MILP. The activity powers $p_a=\mu_a$ are the posterior \emph{mean} of a Bayesian
fit performed \textbf{offline, once, before the loop} by step~0 (\texttt{0\_Regression.jl}): it
reads the raw task-recording data, builds cumulative-$\Delta$SoC energy-balance equations
$A x = b$, and fits $b\sim\mathcal{N}(Ax,s)$ with $\mathrm{TruncatedNormal}(\mu_0,\sigma_0;\ge0)$
power priors and a half-normal noise (pooled NUTS chains). It writes both the mean $\mu$ and a
\textbf{per-activity} standard deviation $\sigma_a$ to \texttt{parameters.csv}. The MILP plans on
$\mu$; the closed-loop plant draws each interval's realized power from
$\mathcal{N}(\mu_a,\sigma_a)$ (truncated at $0$, idle pinned to $0$). The model is \emph{not}
re-fit online during the day.

\end{document}
